\documentclass[10pt]{article}

% Compact course handout. Students submit a NeurIPS-template report in Overleaf;
% this file is the assignment sheet, not the report template.
\usepackage[letterpaper,margin=0.72in]{geometry}
\usepackage[T1]{fontenc}
\usepackage[utf8]{inputenc}
\usepackage{microtype}
\usepackage{amsmath,amssymb}
\usepackage{booktabs,array}
\usepackage[round,authoryear]{natbib}
\usepackage{enumitem}
\usepackage{xcolor}
\usepackage[colorlinks=true,linkcolor=blue!50!black,citecolor=blue!50!black,urlcolor=blue!50!black]{hyperref}

\setlength{\parindent}{0pt}
\setlength{\parskip}{2.5pt}
\setlist[itemize]{leftmargin=1.2em,itemsep=1pt,topsep=2pt,parsep=0pt}
\setlist[enumerate]{leftmargin=1.35em,itemsep=1pt,topsep=2pt,parsep=0pt}
\renewcommand{\arraystretch}{1.08}

\newcommand{\AdamW}{\textsc{AdamW}}
\newcommand{\R}{\mathbb{R}}
\newcommand{\E}{\mathbb{E}}
\newcommand{\repo}{\texttt{[course GitHub repo URL]}}
\newcommand{\overleaf}{\texttt{[editable Overleaf URL]}}

\title{\vspace{-2.2em}\textbf{The Spectral Optimizer Challenge}\\[-0.25em]
\large PhD Optimization for Deep Learning: High Dimensions and Modern GPUs}
\author{}
\date{}

\begin{document}
\maketitle
\vspace{-3.0em}
\thispagestyle{empty}

% ================================================================
\section*{1. Main problem: one-page version}

\textbf{Question.} Can you design an optimizer or optimizer diagnostic that uses \emph{spectral information} from training to improve, explain, or falsify a credible alternative to a carefully tuned \AdamW{} baseline?

\textbf{Core target.} Try to Pareto-dominate tuned \AdamW{} on \emph{both} iterations-to-target and GPU wall-clock on at least one calibrated task, while staying non-catastrophic on the others. A careful negative result is acceptable; an unsupported ``beats Adam everywhere'' claim is not.

\textbf{Choose one spectral signal family.}
\begin{center}
\begin{tabular}{@{}p{0.25\linewidth}p{0.68\linewidth}@{}}
\toprule
\textbf{Curvature spectrum} & Hessian, Fisher/GGN, gradient covariance, Kronecker factors, Hutchinson/Lanczos, saddle curvature. \\
\textbf{Trajectory spectrum} & Recent gradients, momentum, Adam directions, update covariance, Oja/Frequent-Directions sketches, projected energy. \\
\textbf{Spatial spectrum} & Weight/activation/Jacobian singular values, token/channel covariance, residual scaling, dynamical isometry. \\
\bottomrule
\end{tabular}
\end{center}

\textbf{Three tasks.}
\begin{center}
\begin{tabular}{@{}p{0.12\linewidth}p{0.48\linewidth}p{0.28\linewidth}@{}}
\toprule
\textbf{A} & Fashion-MNIST residual MLP & well-conditioned control. \\
\textbf{B} & Synthetic rotated-curvature regression, $\kappa\in\{10,10^3,10^6\}$ & severe anisotropy. \\
\textbf{C} & Teacher--student regression, spectrum slope $\alpha\in\{0.3,1.0\}$ & plateau / slow-spectrum regime. \\
\bottomrule
\end{tabular}
\end{center}

\textbf{Baseline and tuning.} Compare against \texttt{torch.optim.AdamW(foreach=True)} and SGD+momentum. Use the course grid unless your proposal explicitly justifies a smaller sentinel run first:
\[
7\;\text{lr} \times 3\;\text{wd} \times 2\;\epsilon = 42
\]
configs per optimizer per task, followed by three assigned SHA-derived seeds at the best interior configuration.

\textbf{Required final deliverable.} Share \textbf{the editable Overleaf link} using the NeurIPS template. The report abstract must include a clickable GitHub repository link containing code, CSV/JSON artifacts, and a tagged release. \textbf{No paper submission and no separate source/PDF upload.}

\textbf{Non-negotiables.}
\begin{itemize}
  \item No exact \texttt{svd}/\texttt{eigh} on dimension $>256$; no flattening across layers; spectral state is per tensor.
  \item No \texttt{time.time()} for timing claims; use CUDA events and report optimizer-step time separately.
  \item No pre-built Muon/Shampoo/SOAP in the hot path; cite them, compare to them, but do not import them as your method.
  \item No architectural changes to the provided harness unless the project is explicitly a diagnostic, not an optimizer comparison.
\end{itemize}

\vfill
\textbf{What makes a good project?} One falsifiable idea, one correct implementation, one honest comparison, one failure case you understand.

\clearpage
\thispagestyle{empty}

% ================================================================
\section*{2. How to start: one-page setup}

\textbf{First 48 hours.} Do these before inventing anything.
\begin{enumerate}
  \item Clone the harness from \repo{} and run the reference \AdamW{} smoke test.
  \item Implement a disabled mode of your optimizer that is bit- or tolerance-equivalent to \AdamW{}.
  \item Confirm CUDA-event timing: forward, backward, optimizer step, memory peak.
  \item Add one diagnostic log before changing the update rule. Example: top projected energy, effective rank, angle between your update and \AdamW{}.
  \item Write a one-sentence hypothesis: ``My spectral signal helps on [regime] because [mechanism].''
\end{enumerate}

\textbf{Minimal optimizer interface.} Your optimizer should have three modes:
\begin{center}
\begin{tabular}{@{}p{0.24\linewidth}p{0.65\linewidth}@{}}
\toprule
\texttt{disabled} & exactly \AdamW{}; used to prove the harness and parameter groups are correct. \\
\texttt{debug/emit} & computes the proposed update and diagnostics without mutating parameters. \\
\texttt{production} & the actual update used for timed experiments. \\
\bottomrule
\end{tabular}
\end{center}

\textbf{A safe trajectory-style template.} If you choose trajectory spectra, start from a direction-only correction. For a tensor, let $d_t$ be the \AdamW{} direction, $x_t=d_t/\|d_t\|_2$, $B_t\in\R^{r\times d}$ with $B_tB_t^\top=I$, and $c_t=B_tx_t$. Emit
\[
 y_t=\frac{x_t+B_t^\top\Phi_t}{\|x_t+B_t^\top\Phi_t\|_2},
 \qquad \Delta\theta_t=-\eta_t\|d_t\|_2 y_t .
\]
This preserves \AdamW{}'s step magnitude and only changes direction. Log the realized angular perturbation
\[
\tan\theta_t=\frac{\|\Phi_{\perp,t}\|_2}{1+\langle\Phi_t,c_t\rangle},\qquad
\|\Phi_{\perp,t}\|_2^2=\|\Phi_t\|_2^2-\langle\Phi_t,c_t\rangle^2.
\]

\textbf{Pre-register three predictions before the final grid.}
\begin{enumerate}
  \item A regime where your method should win.
  \item A regime where \AdamW{} should win.
  \item A regime where your method should fail, diverge, or become silent.
\end{enumerate}

\textbf{Very brief rubric.} 100 pts: idea/theory 20, implementation 20, empirics 25, ablations/failure analysis 20, clarity/reproducibility 15. Failure analysis is mandatory: a uniform-victory report receives zero in that category.

\textbf{Milestones.}
\begin{center}
\begin{tabular}{@{}p{0.11\linewidth}p{0.78\linewidth}@{}}
\toprule
P1 & One-page proposal, signal family, pseudocode, disabled-\AdamW{} parity. \\
P2 & Working optimizer, diagnostic logs, one sentinel task, first ablation. \\
P3 & Full grid or approved reduced grid, three seeds, failure analysis. \\
P4 & Editable Overleaf NeurIPS report link + GitHub repo link in abstract + oral defense. \\
\bottomrule
\end{tabular}
\end{center}

\vfill
\textbf{Mantra.} If the fast implementation changes the update semantics, it is a different optimizer. Measure before celebrating.

\clearpage

% ================================================================
\appendix
\section*{Appendix A. Reference details}
\addcontentsline{toc}{section}{Appendix A. Reference details}

\subsection*{A.1 What counts as ``beating Adam''}
A claim of Pareto dominance requires: same task, same seeds, same tuning budget, best learning rate not at the boundary, wall-clock measured with CUDA events, and a plot showing both final quality and time-to-target. A result that improves fixed iterations but loses wall-clock should be reported as a fixed-step quality improvement, not an efficiency win. This follows the cautionary optimizer-benchmarking literature: rankings can flip under different hyperparameter spaces or compute budgets \citep{choi2020empirical,teja2020hpo,kaddour2023notrain,algoperf2025}.

\subsection*{A.2 Suggested project tracks}
\begin{tabular}{@{}p{0.24\linewidth}p{0.68\linewidth}@{}}
\toprule
\textbf{Low-rank angular steering} & Implement a small trajectory frame and direction-only correction. Key ablation: correction on/off, rank, update interval, angular distribution. \\
\textbf{Frozen random temporal statistics} & Keep $B_0$ fixed and update only coefficient statistics. Tests whether learning the basis is load-bearing or whether temporal moments suffice. \\
\textbf{Regime-adaptive gates} & Compare static gates against simple regime diagnostics such as effective rank, coefficient coherence, and spectral slope. Must beat the best static mode. \\
\textbf{Implementation shadowing} & Run two implementations on the same frozen stream of directions and compare emitted updates. Useful when a batched GPU path is faster but changes quality. \\
\textbf{Spatial/curvature diagnostic} & Do not change the optimizer first. Measure Jacobian/activation/curvature spectra and predict which task should favor non-diagonal geometry. \\
\bottomrule
\end{tabular}

\subsection*{A.3 Useful diagnostics}
Log at least the following when relevant: update norm, cosine with \AdamW{}, step time, peak memory, active fraction, effective rank, projected energy, $R^2$, signal-to-noise estimate, cap scale, and realized angle $\theta_t$. If a trust region never binds, it is decorative. If a controller emits zero angle, the optimizer is effectively \AdamW{} and the result should say so.

\subsection*{A.4 Common pitfalls}
\begin{itemize}
  \item \textbf{Silent controller:} your diagnostic shows zero activity, but the report claims a mechanism.
  \item \textbf{Boundary-grid win:} best learning rate is the largest/smallest tested value; extend the grid.
  \item \textbf{Mixed CSVs:} multiple SLURM jobs append to one file without job id and metadata.
  \item \textbf{Bad timing:} CPU timers, un-synchronized CUDA, or excluding optimizer overhead.
  \item \textbf{Layer flattening:} a spectral trick works only because it mixes unrelated tensors.
  \item \textbf{Unlabeled figures:} every plot needs task, seeds, optimizer, hyperparameters, units, and whether lower is better.
\end{itemize}

\subsection*{A.5 Why the assignment is strict}
Modern optimizer claims are fragile. \AdamW{} remains a strong baseline; Shampoo/SOAP show that non-diagonal preconditioning can win when engineered and evaluated carefully \citep{shampoo2018,soap2025}; Muon shows that update geometry and parameter routing matter on modern hardware \citep{muon2024,moonlight2025}; first-order norm views explain why changing the update geometry need not be ``second order'' \citep{oldnorm2024}. High-dimensional loss geometry gives motivation, but not a license to overclaim curvature recovery \citep{dauphin2014saddle,pennington2017geometry,saxe2014deeplinear,goldt2019teacherstudent}.

\subsection*{A.6 Deliverable format checklist}
\begin{itemize}
  \item Editable Overleaf link using a NeurIPS template.
  \item Abstract contains a clickable GitHub repository link.
  \item GitHub repo has code, environment file, tagged release, and CSV/JSON artifacts.
  \item Report distinguishes fixed-step quality, wall-clock, and tuning-budget claims.
  \item No paper submission; no separate upload unless the instructor explicitly requests it.
\end{itemize}

\begin{thebibliography}{99}\small
\bibitem[Choi et~al.(2020)]{choi2020empirical}
D. Choi, C. Shallue, Z. Nado, J. Lee, C. Maddison, and G. Dahl.
On empirical comparisons of optimizers for deep learning. 2020.

\bibitem[Teja et~al.(2020)]{teja2020hpo}
P. Teja, F. Mai, T. Vogels, M. Jaggi, and F. Fleuret.
Optimizer benchmarking needs to account for hyperparameter tuning. 2020.

\bibitem[Kaddour et~al.(2023)]{kaddour2023notrain}
J. Kaddour et~al.
No Train No Gain: Revisiting efficient training algorithms for Transformer-based language models. 2023.

\bibitem[Kasimbeg et~al.(2025)]{algoperf2025}
P. Kasimbeg et~al.
Accelerating neural network training: an analysis of the AlgoPerf competition. 2025.

\bibitem[Gupta et~al.(2018)]{shampoo2018}
V. Gupta, T. Koren, and Y. Singer.
Shampoo: preconditioned stochastic tensor optimization. 2018.

\bibitem[Vyas et~al.(2025)]{soap2025}
N. Vyas et~al.
SOAP: Improving and stabilizing Shampoo using Adam. 2025.

\bibitem[Jordan(2024)]{muon2024}
K. Jordan.
Muon: an optimizer for hidden layers in neural networks. 2024.

\bibitem[Liu et~al.(2025)]{moonlight2025}
J. Liu et~al.
Muon is scalable for LLM training. 2025.

\bibitem[Bernstein and Newhouse(2024)]{oldnorm2024}
J. Bernstein and L. Newhouse.
Old Optimizer, New Norm: an anthology. 2024.

\bibitem[Dauphin et~al.(2014)]{dauphin2014saddle}
Y. Dauphin et~al.
Identifying and attacking the saddle point problem in high-dimensional non-convex optimization. 2014.

\bibitem[Pennington and Bahri(2017)]{pennington2017geometry}
J. Pennington and Y. Bahri.
Geometry of neural network loss surfaces via random matrix theory. 2017.

\bibitem[Saxe et~al.(2014)]{saxe2014deeplinear}
A. Saxe, J. McClelland, and S. Ganguli.
Exact solutions to the nonlinear dynamics of learning in deep linear neural networks. 2014.

\bibitem[Goldt et~al.(2019)]{goldt2019teacherstudent}
S. Goldt, M. Advani, A. Saxe, F. Krzakala, and L. Zdeborov\'a.
Dynamics of stochastic gradient descent for two-layer neural networks in the teacher-student setup. 2019.
\end{thebibliography}

\end{document}


